\chapter{봉투에 넣은 예측}

11장의 비교에는 작은 약점이 있어요.
점수도 검사도 같은 여섯 달의 기록에서 나왔다는 거예요.
허락 지도로 만든 점수가 허락을 많이 받은 지갑과 닮은 건 어찌 보면 당연해요.
이미 본 답안지로 채점하는 것과 조금 비슷하지요.

더 어려운 시험이 있어요.
\textbf{미래를 맞히는} 시험이에요.

\section{봉투 실험}

날씨 예보관을 떠올려 보세요.
“내일 비가 온다”고 말한 예보관이 정말 실력이 있는지는 어떻게 알까요?
오늘 예보를 봉투에 넣어 봉해 두고, 내일 진짜 날씨와 비교하면 돼요.
내일 날씨를 보고 나서 예보를 고칠 수 없게 하는 것이 중요해요.

저도 똑같이 했어요.

\begin{enumerate}
\item \textbf{2026년 2월 28일 밤 11시 59분 59초}까지의 기록만으로 지도를 그리고 점수를 매겨요.
      이 점수가 봉투에 든 예보예요. 이날을 \textbf{동결일}이라고 해요.
\item 그다음 \textbf{3월부터 5월까지} 석 달 동안 실제로 일어난 일을 세어요.
\item 예보 줄과 실제 줄을 켄달의 $\tau$로 비교해요.
\end{enumerate}

10장의 달력에서 “13장”이라고 적힌 줄이 바로 이 실험이에요.
봄에 채점했으니 \textbf{봄 홀드아웃}이라고 불러요.
“홀드아웃”은 채점에 쓸 기록을 따로 떼어 놓았다는 뜻이에요.

\section{무엇을 맞히나요?}

이번에 줄을 세운 것은 GMX 거래자들이 아니라, 동결일에 누군가에게서 살아 있는 허락을 받아 둔 통들이에요.
이 통들을 \textbf{스펜더 코호트}라고 불러요.
모두 1,335개였고, 그중 1,174개가 로봇(앱), 161개가 사람의 통이었어요.

맞힐 것은 두 가지예요.

\begin{itemize}
\item \textcolor{allowcol}{\textbf{새 허락}}: 3\textasciitilde{}5월에 새로 생긴 “허락한 통 → 이 통” 쌍이 몇 개인가.
      새 친구가 몇 명이나 문을 열어 주었는가예요. 1,335개 가운데 222개가 새 허락을 받았어요.
\item \textcolor{sendcol}{\textbf{새 송금자}}: 3\textasciitilde{}5월에 처음으로 이 통에 송금한 통이 몇 개인가.
      175개가 새 송금자를 맞았어요.
\end{itemize}

\section{결과}

\begin{figure}[h]
\centering
\begin{tikzpicture}
\begin{axis}[
  ybar, bar width=8pt, width=0.98\linewidth, height=6cm,
  ymin=0, ymax=0.9, ytick={0,0.2,0.4,0.6,0.8},
  symbolic x coords={NA, NS}, xtick=data, xticklabels={새 허락 맞히기, 새 송금자 맞히기},
  enlarge x limits=0.45,
  legend style={font=\scriptsize, at={(0.5,1.02)}, anchor=south, legend columns=3, draw=none},
  nodes near coords, nodes near coords style={font=\tiny, rotate=90, anchor=west},
  every node near coord/.append style={/pgf/number format/fixed, /pgf/number format/fixed zerofill, /pgf/number format/precision=3},
  tick label style={font=\small}, ylabel={$\tau$}, axis x line*=bottom,
]
\addplot[fill=allowcol, draw=allowcol] coordinates {(NA,0.394) (NS,0.283)};
\addplot[fill=sendcol, draw=sendcol] coordinates {(NA,0.123) (NS,0.344)};
\addplot[fill=leafline!70, draw=leafline] coordinates {(NA,0.295) (NS,0.277)};
\addplot[fill=greycol!40, draw=greycol] coordinates {(NA,0.711) (NS,0.404)};
\addplot[fill=greycol!80, draw=greycol] coordinates {(NA,0.117) (NS,0.472)};
\legend{엔도스랭크, 송금 랭크, 동전 랭크, 받은 허락 수, 받은 송금 수}
\end{axis}
\end{tikzpicture}
\caption{봄 홀드아웃(스펜더 1,335개). 막대가 높을수록 미래를 잘 맞힌 거예요.}
\end{figure}

\subsection*{엔도스랭크 대 송금 랭크}

새 허락 맞히기에서 엔도스랭크는 0.394, 송금 랭크는 0.123이었어요.
엔도스랭크가 0.271 앞섰어요.
새 송금자 맞히기에서는 거꾸로였어요.
송금 랭크가 0.344로, 엔도스랭크의 0.283보다 0.061 높았어요.
두 차이 모두 95\% 범위가 0보다 위에 있었어요.

허락 지도로 만든 점수는 새 허락을, 송금 지도로 만든 점수는 새 송금자를 더 잘 맞혔어요.
각자 자기 지도와 닮은 미래를 더 잘 내다본 거예요.
도전자와 챔피언이 한 판씩 나누어 가진 셈이지요.

다만 이 두 비교는 미리 정해 둔 것이 아니에요.
엔도스랭크와 송금 랭크 점수는 봄 채점 결과가 나온 뒤에 매겼거든요.
그래서 저는 숫자를 범위와 함께 그대로 알리기만 하고, 이것으로 무엇을 증명했다고 말하지는 않아요.

\subsection*{송금 랭크가 말할 수 없는 통}

8장에서 엔도스랭크는 허락 지도에 없는 지갑에 대해 아무 말도 할 수 없다고 했지요.
송금 랭크도 마찬가지예요.
스펜더 1,335개 가운데 318개는 동결일까지 다른 통과 송금을 주고받은 기록이 없어서 송금 랭크 점수가 0이었어요.
그런데 이 318개 가운데 22\%가 봄에 새 허락을 받았어요.
나머지 통들은 15\%였는데 말이에요.
송금 랭크는 새 허락을 받을 통이 많은 무리를 줄의 맨 뒤에 세운 셈이에요.

그래서 동결일 전에 송금을 받아 본 977개만 남기고 다시 비교해 보았어요.
그러자 새 허락 맞히기에서 엔도스랭크의 앞섬은 0.271에서 0.042로 줄었어요.
새 송금자 맞히기에서는 송금 랭크가 0.029 앞섰지만, 95\% 범위가 0을 넘나들어서 비슷하다고 봐야 해요.
엔도스랭크가 크게 앞선 데에는, 송금 랭크가 아무 말도 할 수 없는 통이 많았던 탓이 컸던 거예요.
이 점검도 결과를 본 뒤에 한 것이에요.

\subsection*{그런데 더 센 상대가 있었어요}

그래프의 회색 막대를 보세요.
이것은 4장에서 만든, 가장 단순한 점수예요.
동결일까지 받은 화살표를 \emph{세기만} 한 거예요.

받은 허락의 수는 새 허락 맞히기에서 \textbf{0.711}이었어요.
엔도스랭크보다 훨씬 높아요.
받은 송금의 수도 새 송금자 맞히기에서 0.472로, 송금 랭크와 엔도스랭크보다 높았어요.

왜 그럴까요?
이미 허락을 많이 받은 앱은 앞으로도 허락을 많이 받아요.
사람들이 많이 쓰는 앱을 새 사람들도 쓰기 때문이에요.
“부자가 더 부자가 된다”는 말처럼요.
이런 질문에는 복잡한 쪽지 놀이보다 단순한 세기가 더 잘 맞혀요.

저는 이 결과를 논문에 그대로 적었어요.
엔도스랭크도 송금 랭크도, 자기 기록이 앞으로 늘어나는 것을 단순한 세기보다 잘 맞히지는 못했어요.
엔도스랭크가 새 허락에서 송금 랭크를 이긴 건 사실이지만, 그것이 “가장 좋은 방법”이라는 뜻은 아니에요.

\section{동전 랭크는 약속을 지켰을까?}

9장에서 9월 14일에 적어 둔 약속을 기억하나요?
동전 랭크는 새 허락 맞히기에서 앞면 랭크보다 못하지 않고, 새 송금자 맞히기에서 뒷면 랭크보다 나아야 했어요.

\begin{itemize}
\item 새 허락: 동전 랭크 0.295, 앞면 랭크 0.479. 차이가 $-0.184$로 한참 못했어요. \textcolor{roseline}{\textbf{×}}
\item 새 송금자: 동전 랭크 0.277, 뒷면 랭크 0.270. 차이가 $+0.007$인데 95\% 범위가 $-0.012$에서 $0.028$로 0을 넘나들어요. 낫다고 할 수 없어요. \textcolor{roseline}{\textbf{×}}
\end{itemize}

두 조건 모두 넘지 못했어요.
동전 랭크는 9월 14일의 약속을 지키지 못했어요.
$\lambda$를 0.25나 0.75로 바꿔도 마찬가지였어요.

\section{그런데 눈에 띈 것이 하나 있었어요}

결과를 들여다보다가 약속에 없던 것이 눈에 띄었어요.
동전 랭크를 뒷면 랭크하고만 비교하면, 새 허락 맞히기에서 동전 랭크가 0.295로 뒷면 랭크의 0.044보다 \textbf{0.251}이나 높았어요.
새 송금자 맞히기에서는 뒷면 랭크와 비슷했고요.

“송금 지도에 허락 화살표를 섞으면 송금 지도만 걷는 것보다 나아진다.”
이것은 사실 제 논문 제목이 말하고 싶은 바로 그 이야기예요.
하지만 저는 이것을 결과라고 부를 수 없었어요.
\emph{결과를 보고 나서 찾은 것}이기 때문이에요.
송금 랭크와 견주어 보면 더 조심해야 해요.
새 허락에서는 동전 랭크가 앞섰지만(0.295 대 0.123), 새 송금자에서는 뒤졌거든요(0.277 대 0.344).
왜 결과를 보고 나서 찾은 것이 문제인지, 그래서 제가 무엇을 했는지가 다음 장의 이야기예요.

\begin{deeper}[진짜 숫자]
\begin{center}
\small
\begin{tabular}{@{}lcc@{}}
\toprule
점수 (2월 28일 동결) & 새 허락 $\tau$ & 새 송금자 $\tau$ \\
\midrule
엔도스랭크 & 0.394 & 0.283 \\
송금 랭크 & 0.123 & 0.344 \\
앞면 랭크 (동전 랭크, $\lambda=1$) & 0.479 & 0.372 \\
뒷면 랭크 (동전 랭크, $\lambda=0$) & 0.044 & 0.270 \\
동전 랭크 ($\lambda=0.5$) & 0.295 & 0.277 \\
씨앗 랭크 & 0.114 & 0.264 \\
받은 허락 수 & 0.711 & 0.404 \\
받은 송금 수 & 0.117 & 0.472 \\
\bottomrule
\end{tabular}
\end{center}
엔도스랭크 $-$ 송금 랭크는 새 허락 $+0.271$ [0.207, 0.332], 송금 랭크 $-$ 엔도스랭크는 새 송금자 $+0.061$ [0.026, 0.098]이에요(둘 다 결과를 본 뒤 계산).
송금을 받아 본 스펜더 977개에서는 엔도스랭크 $-$ 송금 랭크가 새 허락 $+0.042$ [0.003, 0.085], 새 송금자 $-0.029$ [$-0.070$, 0.011]이었어요.
허락한 양을 그대로 쓰는 앞면 랭크는 봄에 두 시험 모두 엔도스랭크보다 높았어요.
스펜더는 블록체인 노드에 코드가 있는지 물어서 계약 1,174개, 사람의 지갑 161개(계약 기능을 덧붙인 새 방식의 지갑 1개 포함)로 나누었어요.
\end{deeper}

\begin{learned}
\begin{itemize}
\item 2월 28일에 점수를 봉해 두고, 3\textasciitilde{}5월에 실제로 일어난 일로 채점했어요.
\item 새 허락은 엔도스랭크가, 새 송금자는 송금 랭크가 더 잘 맞혔어요. 각자 자기 지도와 닮은 미래를 잘 맞힌 거예요.
\item 엔도스랭크가 크게 앞선 데에는 송금 랭크가 아무 말도 할 수 없는 통이 많았던 탓이 커요.
\item 받은 화살표를 세기만 한 단순한 점수가 둘 다보다 잘 맞혔어요.
\item 동전 랭크는 9월 14일의 약속을 지키지 못했어요.
\item 동전 랭크가 뒷면 랭크보다 새 허락을 잘 맞힌 것은 결과를 보고 나서 찾은 것이라, 아직 결과라고 부를 수 없었어요.
\end{itemize}
\end{learned}
